\section{Conclusion}
\label{sec:conclusion}
We introduce \textbf{AgentJudgeBench}, a benchmark for measuring LLM-judge reliability on structured, dependency-driven tool-calling - a setting where existing LLM-as-judge work offers little calibration. Three design choices distinguish it: difficulty-stratified rewrites (easy / medium / hard preserving the ground-truth trace), per-record DAG-topology annotation across six patterns, and a paired with-GT / no-GT evaluation protocol.

Across $321{,}648$ completed evaluations, four patterns emerge. (1)~Alignment degrades monotonically with difficulty, approximately $1.5\times$ faster without ground truth ($-12$ to $-14$\,pp) than with it ($-6$ to $-11$\,pp). (2)~Hard no-GT alignment converges to a $77$-$82\%$ band regardless of judge capacity - a task-level ceiling confirmed by the C2 recalibrated-prompt ablation (${\leq}{+1.0}$\,pp for capable generators). (3)~GT exposure is counterproductive for frontier judges (Gemini-2.5-Pro: $-3.9$\,pp; GPT-5.4: $-1.5$\,pp), consistent with over-anchoring. (4)~Prompt structure is the dominant configuration lever: structured rubrics add $+4.8$-$+6.5$\,pp; CoT reasoning adds ${\leq}0.3$\,pp. The practical guidance is direct: \textbf{with ground truth, use QwQ-32B; without it, frontier judges lead narrowly but the no-GT ceiling limits the choice.} Limitations and deployment guidance are in Appendices~\ref{app:limitations} and~\ref{app:decision-guide}.